\documentclass[11pt]{article}
\usepackage{fontspec,amsmath,amssymb,amsthm,hyperref}
\usepackage[margin=1in]{geometry}
\setlength{\emergencystretch}{3em}
\title{Actual splitting laws, L\ensuremath{{}^{2}} realizations and iid law transfer}
\author{LangevinLean mathematical exposition}
\date{}
\begin{document}\maketitle
\noindent Source module: \texttt{ProximalLangevin/SplittingLaw.lean}. Compile this proof document with XeLaTeX or LuaLaTeX. The accompanying HTML pairs the mathematical argument with the exact Lean statements and records the scope of the formalization.


The pointwise splitting update applies the averaged true proximal map after a noisy smooth step. Convexity and a lower bound for \(g_2\), positive \(\lambda\), \(N>0\) and a Lipschitz smooth gradient make this update measurable and preserve finite quadratic moments.


\section{Prove measurability of the exact pointwise update}


Nonexpansiveness proves \(\operatorname{Lip}(\operatorname{Pair}_{P_\lambda})\le1\). The forward map \((z,\xi)\mapsto z-(\gamma/N)\nabla F_{g_1}(z)+\sqrt{2\gamma/N}\xi\) is measurable by continuity of its terms; compose it with proximity. Induction proves equality of the custom pointwise chain with the generic noise-driven chain.




\textbf{Correspondence with the displayed declarations.}


\begin{itemize}

\item \texttt{pipglaUpdate\_measurable}: The exact splitting update is jointly measurable in the state and noise.
\item \texttt{pipglaChain\_eq\_noiseDrivenChain}: The pointwise recursion is the existing generic noise-driven chain.
\end{itemize}


\noindent Lean declarations: \texttt{pairwiseMap\_prox\_lipschitz}, \texttt{pipglaUpdate\_measurable}, \texttt{pipglaChain\_eq\_noiseDrivenChain}.

\section{Realize the actual chain inside L\ensuremath{{}^{2}}}


Define the algorithm law as the pushforward by the explicit chain. Define an \(L^2\) update by applying the Lipschitz pairwise map to the \(L^2\) noisy Euler update. Their representatives agree almost everywhere by the composition and addition identities in \(L^2\). Inductively every iterate lies in \(L^2\) and agrees almost everywhere with its pointwise recursion. This constructs moment control without Gaussianity or equilibrium assumptions.




The square-integrable update first applies the Euler drift/noise map and then the genuine pairwise proximal map. Both are Lipschitz, so the result is again an \(L^2\) random variable. The update representative theorem intersects the full-measure sets for these two operations and obtains the pointwise PIPGLA formula. An induction repeats this at each fixed finite time, replacing the previous representative by the pointwise iterate on their common full-measure event. The chain's pushforward law is therefore exactly \texttt{pipglaLaws}, rather than merely a convenient auxiliary recursion.




\textbf{Correspondence with the displayed declarations.}


\begin{itemize}

\item \texttt{pipglaLaws}: The law of the actual PIPGLA state after a supplied innovation sequence.
\item \texttt{lpPipglaUpdate}: The actual noisy forward step followed by the genuine pairwise proximal map, represented in L\ensuremath{{}^{2}}.
\item \texttt{lpPipglaChain}: Every iterate is an actual L\ensuremath{{}^{2}} random variable.
\end{itemize}


\noindent Lean declarations: \texttt{pipglaLaws}, \texttt{lpPipglaUpdate}, \texttt{lpPipglaUpdate\_coe}, \texttt{lpPipglaChain}, \texttt{lpPipglaChain\_coe}.

\section{Pushforward equality proves mass one}


Almost everywhere equal representatives have equal pushforward measures, so \(\rho_n=Q\circ X_n^{-1}\) for the constructed \(L^2\) iterate. Its almost everywhere measurability makes this a genuine probability law: the preimage of the whole state space has \(Q\)-mass one.



\noindent Lean declarations: \texttt{pipglaLaws\_eq\_lpLaw}, \texttt{pipglaLaws\_isProbability}.

\section{Transfer any iid realization to the canonical product law}


Assume the initial variable is independent of the whole innovation path and the innovations form an independent family with common law \(\eta\). The joint input law is \(\mathcal L(X_0)\otimes\eta^{\otimes\mathbb N}\). Applying the measurable noise-driven recursion gives exactly the canonical chain law. The proof uses the actual independence and common-law hypotheses, without assuming a stochastic error estimate.




For a measurable update, the time-\(n\) state is a measurable function of initialization and the innovation path. Independence of initialization from that path makes their joint law a product; iid innovations make the path law \(\eta^{\otimes\mathbb N}\). Pushing this product through the finite recursion proves the canonical-law equality. The assumptions \texttt{hinit}, \texttt{hnoise}, \texttt{hind}, \texttt{hiid} and \texttt{hlaw} respectively supply these measurability, independence and distribution facts; none is a convergence-rate premise.




\textbf{Correspondence with the displayed declarations.}


\begin{itemize}

\item \texttt{pipglaLaws\_eq\_canonical}: Any independent iid realization has the canonical splitting law. This proves law transfer across coupling spaces without a convergence assumption.
\end{itemize}


\noindent Lean declarations: \texttt{pipglaLaws\_eq\_canonical}.

\end{document}
